% Wampanoag (Winslow 1624) — Source Workflow
% Issue #1366
% Status: Research complete, implementation pending
% !TEX program = xelatex
% !TEX encoding = UTF-8 Unicode

\section{Wampanoag: Winslow 1624}
\label{sec:1366}

\subsection{Source Description}

\textbf{Recorder:} Edward Winslow (1595--1655), Mayflower passenger, Plymouth
Colony governor for three terms (1633--1634, 1636--1637, 1644--1645), and the
colony's primary negotiator and diplomat with the Native peoples of New England.

\textbf{Source:} \emph{Good Newes from New England} (1624), printed in London by
I.D.\ for William Bladen and John Bellamie. The earliest published source of
Wampanoag vocabulary in the SNEA corpus.

\textbf{Language:} Wampanoag, Pokanoket dialect [wam].

\textbf{Corpus:} 24 records (the smallest in the SNEA collection).

\textbf{Orthography:} English-hearer perception with zero regularization.
Winslow was unusually honest about his linguistic limitations, describing the
Massachusett language as ``very copious, large, and difficult. As yet we cannot
attain to any great measure thereof.'' This acknowledgment distinguishes
Winslow from many other colonial recorders who pretended to greater competence.
His transcriptions are therefore likely more reliable as raw phonetic
impressions --- he was not normalizing or regularizing what he heard.

\subsection{Recorder Biography}

Edward Winslow was among the 41 Mayflower passengers who signed the Mayflower
Compact in November 1620. Key events relevant to his linguistic data:

\begin{itemize}
  \item \textbf{March 1621:} First contact with Samoset (Abenaki sachem speaking
    English) and Massasoit.
  \item \textbf{March--July 1621:} Tisquantum (Squanto) serves as interpreter,
    teaching the Pilgrims Wampanoag vocabulary.
  \item \textbf{March 1623:} Winslow cures Massasoit of a life-threatening
    illness --- the encounter that produced his most extensive Wampanoag
    dialogue.
  \item \textbf{August 1623:} Hobbamock warns Plymouth of Massasoit's alleged
    conspiracy with the Massachusett; Standish's Wessagussett raid follows.
  \item \textbf{1624:} \emph{Good News from New England} published in London.
\end{itemize}

After returning to England in 1646, Winslow served as Oliver Cromwell's chief
agent in foreign missions. He died of yellow fever in 1655 en route to a
diplomatic mission in Spanish Hispaniola.

\subsection{Linguistic Intermediaries}

Winslow did not learn Wampanoag directly --- he worked through interpreters.
The identity and linguistic background of these intermediaries is crucial for
understanding potential biases in his transcriptions.

\subsubsection{Tisquantum (Squanto, c.~1585--1622)}

Tisquantum was a Patuxet Wampanoag with an extraordinary linguistic history:

\begin{itemize}
  \item \textbf{1614:} Kidnaped by Captain Thomas Hunt, along with 20+ other
    Patuxet men.
  \item \textbf{1614--1618:} Sold as a slave in Málaga, Spain; escaped with the
    help of Spanish friars.
  \item \textbf{1618--1619:} Lived in London with John Slany, a merchant of the
    Newfoundland Company; learned English.
  \item \textbf{1619:} Returned to Patuxet (modern Plymouth) --- only to find
    his entire village depopulated by the 1616--1619 epidemic.
  \item \textbf{1621:} Approached Plymouth Colony in March; served as
    interpreter and guide.
  \item \textbf{1622:} Died of ``Indian fever'' at Monomoit (Chatham, Cape Cod)
    while accompanying Winslow on a diplomatic mission.
\end{itemize}

Linguistic implications of Tisquantum's mediation:
\begin{itemize}
  \item His English was learned in London, not from the Pilgrims --- it was
    16th-century English, not colonial dialect.
  \item His Wampanoag was Patuxet, a sub-dialect of Wampanoag closely related
    to Pokanoket but not identical.
  \item His time in Spain and England (8+ years) may have affected his native
    language production (attrition).
\end{itemize}

\subsubsection{Hobbamock (Hobomok, d.~1642)}

Hobbamock was a Pokanoket \emph{pniese} --- a military champion and elite
warrior --- who served as Massasoit's trusted advisor and was assigned to live
among the English:

\begin{itemize}
  \item \textbf{March 1621:} Assigned by Massasoit to reside at Plymouth Colony
    as a liaison.
  \item \textbf{March 1623:} Revealed Massasoit's alleged conspiracy with the
    Massachusett, leading to Standish's raid.
  \item \textbf{Accompanied Winslow} on the mission to cure Massasoit --- his
    lament ``Neen womasu sagimus'' (``My loving sachem'') upon hearing of
    Massasoit's apparent death is one of Winslow's most famous transcriptions.
  \item \textbf{Died of European disease} in 1642.
\end{itemize}

Linguistic significance: Hobbamock likely spoke Pokanoket dialect natively.
He was not a language teacher but may have provided a more natural sample of
Wampanoag speech than Tisquantum. His English was likely limited --- Winslow
probably communicated with him through a mixture of Wampanoag and signs.

\subsection{Complete Vocabulary}

According to Salvucci (InteractiveALR Vol.\ 27) and confirmed by primary source
reading, Winslow's text contains 22 distinct words and phrases in Wampanoag:

\begin{table}[htbp]
\centering
\caption{Winslow's Wampanoag Lexicon}
\label{tab:1366-lexicon}
\begin{tabular}{lll}
\hline
Word/Phrase & English & Notes \\
\hline
Keen Winsnow? & ``Art thou Winslow?'' & Massasoit's greeting; l$\rightarrow$n \\
Ahhe & ``Yes'' & Winslow's reply \\
Matta neen wonckanet namen Winsnow & ``O Winslow, I shall never see thee again!'' & Deathbed lament \\
Winsnow & ``Winslow'' & /l/ $\rightarrow$ /n/; no /l/ in Wampanoag \\
Neen womasu sagimus & ``My loving sachem'' & Hobbamock's lament \\
Sachimo comaco & ``The sachem's house'' & Sachem + house at Pokanoket \\
Witeo & ``An ordinary house'' & Contrasts with comaco \\
Squasachim & ``The sachem's wife'' & Squa (woman) + sachim (chief) \\
Maskiet & ``Physic/Medicine'' & Healing context \\
Kiehtan & (The supreme God) & Explains Wampanoag theology \\
Chise & ``An old man'' & Winslow's etymology for Kiehtan \\
Kiehchise & ``A very old man'' & Comparative: chise $\rightarrow$ kiehchise \\
Hinnaim namen, hinnaim michen, matta cuts & ``By and by see, eat, but not speak'' & Complex sentence \\
Quatchet & ``Walk abroad'' & Imperative? \\
Sachim & ``King/Chief'' & General title \\
Pniese (also pinse, pinese) & ``Chief champion / man of valor'' & Military elite class \\
Powah & ``Shaman / healer'' & Calls upon the devil \\
Squaw & ``Woman'' & Occurring in compounds \\
Massasoit & ``Great Chief'' & Proper name / title \\
Canonickus & Massasoit's older brother & Proper name with -us? \\
Wamsutta & Massasoit's son & Proper name \\
Namasquemit & ``Powah'' or dreaming healer & One reference only \\
\hline
\end{tabular}
\end{table}

\subsection{Syntactic Complexity}

Despite the small corpus, Winslow records several sentences showing Wampanoag
syntax:

\begin{description}
  \item[Keen Winsnow?] Copular question: \emph{Keen} = 2sg pronoun ``thur'',
    predicate nominal.
  \item[Matta neen wonckanet namen Winsnow] Negative declarative: \emph{Matta}
    (neg) + \emph{neen} (1sg) + \emph{wonckanet} (again/future) + \emph{namen}
    (see) + vocative.
  \item[Hinnaim namen, hinnaim michen, matta cuts] Serial verb construction:
    three coordinate clauses with \emph{hinnaim} (by-and-by) + \emph{matta}
    (neg).
  \item[Neen womasu sagimus] Possessive + noun: \emph{neen} (I/my) +
    \emph{womasu} (loving?) + \emph{sagimus} (sachem + 1sg possessive).
\end{description}

\subsection{Orthography System}

Winslow wrote in Early Modern English orthography (pre-1625), before English
spelling standardization. Understanding what Winslow's spellings meant for
English is essential before mapping to Wampanoag phonemes.

\subsubsection{Vowel Grapheme-to-Phoneme Mapping}

\begin{table}[htbp]
\centering
\caption{Winslow Vowel Mapping}
\label{tab:1366-vowels}
\begin{tabular}{llll}
\hline
Grapheme & IPA & Example & Confidence \\
\hline
a (short) & /a/ $\sim$ /ə/ & matta & High \\
a (long) & /aː/ & ahhe (first a) & Moderate \\
e (short) & /ɛ/ & chise (final e) & Moderate \\
e (final) & /ə/ & ahhe, witeo, maskiet & High \\
i & /ɪ/ $\sim$ /i/ & kiehchise, chise & Moderate \\
o & /ɔ/ $\sim$ /ə/ & hobbamock, commaco & Moderate \\
u & /ʌ/ $\sim$ /ə/ & cuts & Low \\
ie & /iː/ & kiehtan, kiehchise & High \\
oo & /uː/ & askooke & High \\
ah & /ah/ & ahhe & Moderate \\
ee & /iː/ & keen (thou) & High \\
\hline
\end{tabular}
\end{table}

\subsubsection{Consonant Grapheme-to-Phoneme Mapping}

\begin{table}[htbp]
\centering
\caption{Winslow Consonant Mapping}
\label{tab:1366-consonants}
\begin{tabular}{llll}
\hline
Grapheme & IPA & Example & Confidence \\
\hline
h & /h/ & ahhe, hobbamock, kiehtan & High \\
k & /k/ & kiehtan, askooke, keen & High \\
c (before a/o) & /k/ & commaco, cuts, comaco & High \\
ch & /tʃ/ & chise, kiehchise & High \\
sh & /ʃ/ & Only in proper names & Low \\
m & /m/ & commaco, matta, namen & High \\
n & /n/ & neen, namen, Winsnow & High \\
s & /s/ & sachim, squaw, sagimus & High \\
t & /t/ & matta, kiehtan, quatchet & High \\
w & /w/ & winsnow, witeo & High \\
p & /p/ & pniese (initial pn-!) & Moderate \\
b & /b/ & hobbamock (medial) & Moderate \\
\hline
\end{tabular}
\end{table}

\subsubsection{Preaspiration /h/ in Winslow}

Winslow writes /h/ more consistently than any other early Wampanoag source
except Edwards (1787 Mahican). This is one of the most valuable aspects of his
small corpus.

\begin{table}[htbp]
\centering
\caption{Cross-Source Comparison for /h/}
\label{tab:1366-h}
\begin{tabular}{llll}
\hline
Source & hobbamock initial h & kiehtan medial ht & Overall /h/ Reliability \\
\hline
Winslow (1624) & YES & YES & Highest among early sources \\
Wood (1634) & NO (abamocho) & N/A & Unreliable for /h/ \\
Williams (1643) & Mixed (ábammocho) & N/A & Unreliable \\
Eliot (1663) & Different root & Likely kehtan & Partial \\
Edwards (1787) & N/A (Mahican) & YES & Most reliable overall \\
\hline
\end{tabular}
\end{table}

Interpretation: Winslow's relatively good /h/ transcription may reflect:
\begin{itemize}
  \item His intermediaries (Tisquantum, Hobbamock) had clearer /h/ articulation
  \item The Pokanoket dialect had more salient /h/ than Natick
  \item Winslow was a more careful listener --- his ``ignorance'' of the
    language made him more phonetically accurate (less normalization bias)
\end{itemize}

\subsection{Pidginization Analysis (Goddard 1977)}

Goddard's landmark analysis of early contact Algonquian data identified
Winslow's vocabulary as showing clear evidence of pidginization --- the English
and Wampanoag were communicating via a simplified contact variety, not the full
grammatical language.

\subsubsection{Key Evidence}

In the sentence \emph{Matta neen wonckanet namen Winsnow}:

\begin{enumerate}
  \item \textbf{\emph{namen} (transitive inanimate ``see it'')} is used instead
    of the grammatically correct transitive animate form (expected \emph{nunaw}
    ``I see him''). The verb agrees with the object for gender --- \emph{namen}
    treats ``Winslow'' as if he were an inanimate object.
  \item \textbf{Pro-drop is lost:} Full independent pronouns (\emph{neen} ``I'',
    \emph{keen} ``thou'') are used where verbal inflection alone would carry
    person in the full language. This is a universal pidgin feature ---
    morphology simplifies, syntax becomes analytic.
  \item \textbf{No verbal inflection:} The expected inflected verb
    (\emph{n-nam-en} = 1sg-see-TA) is replaced with what appears to be the
    citation form or a nominalized form.
\end{enumerate}

\subsubsection{The Massachusett Pidgin Hypothesis}

Goddard hypothesized that a Massachusett Pidgin (similar in function to Delaware
Pidgin or Mobilian Jargon) existed before English colonization, used for
inter-tribal trade. If so, what Winslow recorded may not have been ``pure''
Wampanoag at all, but rather a trade pidgin. This would explain:

\begin{itemize}
  \item The consistency of certain simplifications (pidgin features shared
    across sources)
  \item The spread of certain forms (e.g., \emph{matta} for negation appears in
    all colonial sources)
  \item The divergence between ``missionary'' Massachusett (Eliot 1663) and
    ``contact'' Massachusett (Winslow, Wood)
\end{itemize}

\subsection{Cross-Source Cognate Table}

\begin{table}[htbp]
\centering
\caption{Confirmed Cognates Across Sources}
\label{tab:1366-cognates}
\begin{tabular}{lllll}
\hline
Gloss & Winslow (1624) & Wood (1634) & Trumbull/Eliot & Williams (1643) \\
\hline
`no' & matta & matta & matta & matta \\
`woman' & squa(achim) & squaw & squa & squa \\
`devil' & hobbamock & abamocho & achm∞wonk & ábammocho \\
`yes' & ahhe & --- & --- & --- \\
`house' & witeo & wetu/weetu & wetu & --- \\
`snake' & askooke & askook & askkuhnk & askug \\
`king' & sachim & sachim & sachim & --- \\
`old man' & chise & --- & --- & --- \\
`speak' & cuts & --- & cuts & ketes \\
`walk' & quatchet & --- & --- & --- \\
`medicine' & maskiet & --- & maskiht & --- \\
`shaman' & powah & powaw & --- & powwáw \\
\hline
\end{tabular}
\end{table}

\subsection{Dialect Affiliation: Pokanoket Wampanoag}

Winslow's data was collected primarily at Plymouth (Patuxet) and Pokanoket
(modern Bristol, RI --- Massasoit's capital). The Wampanoag Confederacy in 1620
consisted of several bands: Pokanoket (Massasoit's people, primary), Patuxet
(Plymouth, Tisquantum's people, depopulated by 1620), Nauset (Cape Cod),
Monomoit (Chatham), Manomet (Sandwich/Bourne), and Mashpee (Cape Cod).

The key dialectal division within Massachusett/Wampanoag:

\begin{table}[htbp]
\centering
\caption{Pokanoket vs.\ Natick Features}
\label{tab:1366-dialect}
\begin{tabular}{llll}
\hline
Feature & Pokanoket (Winslow/Wood) & Natick (Eliot/Trumbull) & Notes \\
\hline
`house' & witeo (Winslow) & wetu (Eliot) & Real lexical difference? \\
`devil' & hobbamock (Winslow) & achm∞wonk (Eliot) & Different roots entirely \\
`man of valor' & pniese (Winslow) & Not attested & Nativized English? \\
Initial /h/ & Written (Winslow) & Often omitted (Eliot) & Possibly orthographic \\
/l/ substitution & n only (Winslow) & n only (Eliot) & Consistent across dialects \\
\hline
\end{tabular}
\end{table}

Per Salvucci's assessment: ``Winslow's data has not been systematically
analyzed for dialect affiliation.'' The Pokanoket sub-corpus is so small (22
items) that confident dialect assignment is impossible without external
evidence.

\subsection{Problems Encountered}

\begin{itemize}
  \item \textbf{Tiny corpus:} 24 records is insufficient for any statistical
    approach. Every entry must be treated as a unique case. Winslow is a
    corroborative source only.
  \item \textbf{Pidginization:} Goddard (1977) identifies ``a good deal of
    pidginization'' in Winslow's data. This affects whether the data can be
    used for phonological reconstruction at all. If the English were learning a
    pidgin variety, word order may not reflect full Wampanoag syntax, verb forms
    may be bare stems, and vocabulary may be mixed from multiple Algonquian
    languages.
  \item \textbf{Intermediary filter:} Tisquantum (Squanto) and Hobbamock served
    as interpreters --- their own linguistic backgrounds may have affected the
    transcriptions. Tisquantum's English was learned in London, not from the
    Pilgrims, and his Wampanoag was Patuxet, a sub-dialect closely related to
    but not identical with Pokanoket.
  \item \textbf{Dialect uncertainty:} Pokanoket vs.\ Natick differences are
    unresolved. The distinction may represent real dialect differences, register
    differences (pidgin vs.\ full language), chronological differences (1620s
    vs.\ 1660s), or orthographic convergence.
  \item \textbf{Proper name contamination:} Some entries include English names
    (Winsnow = Winslow), which follow English orthography and must be excluded
    from phonemic analysis.
  \item \textbf{Multi-word entries:} Several entries are entire sentences
    requiring syntactic segmentation before lexical analysis.
  \item \textbf{Uncertain dialect:} Pokanoket dialect vs.\ Patuxet (via
    Tisquantum) vs.\ generic pidgin. Not resolvable from Winslow's data alone.
\end{itemize}

\subsection{Research Approach}

Five-step corroboration pipeline using Trumbull and Wood as primary sources.
Winslow data is treated as corroborative, not primary. Each entry is verified
against expected Eliot/Trumbull forms.

\subsubsection{Step 1: Corroboration-Only}

Treat Winslow as corroborative evidence for rules derived from larger sources.
For each Winslow entry, find the cognate in larger Wampanoag sources (Trumbull
4,491 records, Wood 339 records). Use the cross-source consensus for the
phonemic mapping. Flag Winslow-only words for manual review.

\subsubsection{Step 2: Key Cognate Mappings}

\begin{table}[htbp]
\centering
\caption{Key Cognate Mappings}
\label{tab:1366-key-cognates}
\begin{tabular}{llll}
\hline
Winslow & Wood & Trumbull/Eliot (expected) & PA Source \\
\hline
hobbamock (devil) & abamocho & Different root & See §11 --- two different roots? \\
matta (no) & matta & matta & PA *mati- \\
kiehtan (God) & --- & kehtan & PA *ki·ht-? \\
sachim (chief) & sachim & sachim & PA *sa·kima·wa \\
squa (woman) & squaw & squa & PA *e·θkwe·wa \\
askooke (snake) & askook & askkuhnk & PA *aθkw-? \\
\hline
\end{tabular}
\end{table}

\subsubsection{Step 3: h-Transcription as Secondary Evidence}

Winslow's relatively consistent /h/ transcriptions provide secondary evidence
for where preaspiration occurs:
\begin{itemize}
  \item When Winslow writes \emph{h} and other sources also confirm it
    $\rightarrow$ well-attested
  \item When Winslow writes \emph{h} and other sources omit it $\rightarrow$
    possible but needs external confirmation
  \item When Winslow omits \emph{h} and Edwards (Mahican) writes it
    $\rightarrow$ probable preaspiration unrecorded
\end{itemize}

\subsubsection{Step 4: Segment k- Prefix}

The 2nd person prefix \emph{k-} is seen in \emph{keen} (thou, 2sg pronoun) and
should be segmented before conversion. Note that \emph{k-} with vowel-initial
stems would produce different surface forms than \emph{k-} with consonant-initial
stems.

\subsubsection{Step 5: Verify Against Trumbull}

For each converted Winslow entry, verify against the expected form in the
Eliot/Trumbull orthography. Flag mismatches for manual review --- especially:
\begin{itemize}
  \item Where Winslow writes /h/ and Trumbull does not (possible preaspiration
    evidence)
  \item Where Winslow and Trumbull use different lexical roots (possible dialect
    divergence)
  \item Where Winslow's form contradicts known PA correspondences
\end{itemize}

\subsection{Conversion Workflow}
\texttt{[Pending implementation --- see Issue \#1366]}

The conversion pipeline will be implemented as a Python module in
\texttt{src/commons/parsing/}. The module will:

\begin{itemize}
  \item Treat Winslow as corroborative only --- no regex-based conversion
  \item For each Winslow entry, find the cognate in larger Wampanoag sources
  \item Use cross-source consensus for phonemic mapping
  \item Flag Winslow-only words for manual review
  \item Segment \emph{k-} prefix before conversion
  \item Verify each converted entry against expected Eliot/Trumbull forms
\end{itemize}

\subsection{Linguist Feedback Template}

\texttt{[Template --- content pending review]}

The following questions are open for linguist review:

\begin{enumerate}
  \item \textbf{Pidgin or genuine dialect?} Are the simplifications in Winslow's
    data (e.g., \emph{namen} for ``see'' used as transitive animate) evidence
    of a pre-existing Massachusett Pidgin (Goddard 1977) or simply a learner's
    imperfect recording? This affects whether Winslow's data can be used for
    phonological reconstruction at all.
  \item \textbf{\emph{witeo} vs.\ \emph{wetu}:} Is the \emph{witeo} (Winslow) /
    \emph{wetu} (Eliot) split evidence of Pokanoket vs.\ Natick dialect
    difference, or are both from PA *wi·kiwa·mi with different English
    ear-perceptions? Williams records Narragansett \emph{wetu} (agreeing with
    Eliot), suggesting \emph{witeo} may be uniquely Pokanoket.
  \item \textbf{\emph{hobbamock} volcanism:} Is \emph{hobbamock} a genuine
    Wampanoag lexical item or a portmanteau of Wampanoag + English
    ``hobgoblin''? The initial \emph{h} and the \emph{-ock} suffix appear in
    both the proper name and the common noun ``devil'', suggesting folk
    etymology contamination is possible.
  \item \textbf{Phonological reliability of pidgin data:} If pidgin phonology
    approximates the substrate, then Winslow's transcriptions are phonologically
    useful. If the pidgin introduced phonological simplifications (cluster
    reduction, final devoicing, vowel merger), then even the phonology is
    unreliable.
  \item \textbf{\emph{pniese} etymology:} Is this the same root as PA
    *-epye·wa (``to be powerful, to be a man'')? The initial cluster /pn-/ is
    unusual for Massachusett and may reflect a specific morphotactic condition
    or epenthesis.
\end{enumerate}

\subsection{Related Work: Colonial Recorder Perception Problem}

The core challenge of this research card --- English-speaking recorders
misperceiving and mis-transcribing Algonquian phonemes --- is a known
phenomenon documented across Eastern North America. Two researchers' work is
directly relevant:

\textbf{Dr.\ Keith Cunningham} (Linguist, Nanticoke Indian Tribe). His doctoral
dissertation \emph{A Phonological Analysis of Nanticoke With Practical
Applications for Language Revitalization} (Georgetown University) analyzes three
18th-century Nanticoke word lists (Heckewelder 1785) using the same methodology:
historical phonology from colonial recorders, with the same challenges of
multilingual informants and orthographic complexity. His 2024 Algonquian
Conference presentation \emph{A Revised Phonological Analysis of the Heckewelder
Vocabulary of Nanticoke} directly parallels the Winslow analysis --- both
involve tiny corpora (Winslow 24 records, Heckewelder \textasciitilde200)
recorded by English speakers with no linguistic training, both require PA
comparative calibration, and both must account for intermediary effects
(Tisquantum/Hobbamock for Winslow, multilingual informants for Heckewelder).
Cunningham's work demonstrates that the colonial recorder perception problem
extends beyond SNEA into the broader Eastern Algonquian family.

\textbf{Dr.\ Craig Kopris} (Independent scholar). His paper \emph{Les sons
wyandots perçus par des oreilles étrangères} (Wyandot sounds perceived by
foreign ears, \emph{Recherches amérindiennes au Québec}, 2014) is the exact
framing of the Winslow problem: how European recorders systematically
misperceived and mis-transcribed indigenous phonemes. His \emph{Wyandot
Phonology: Recovering the Sound System of an Extinct Language} applies the same
recovery methodology to an unrelated language family (Iroquoian), demonstrating
that the foreign-ear perception problem is universal across colonial
documentation of North American languages.

The Winslow corpus (24 records, the smallest in the SNEA collection) represents
the most extreme case of the tiny-corpus problem. The methodological question
--- whether such a small dataset can yield reliable phonological data --- is
directly addressed by both Cunningham's work on small Nanticoke word lists and
Kopris's work on recovering phonology from limited documentary sources.

\subsection{Key References}

\begin{itemize}
  \item Aubin, George F. 1975. \emph{A Proto-Algonquian Dictionary}. Ottawa:
    National Museums of Canada.
  \item Goddard, Ives. 1977. ``Massachusett Pidgin.'' \emph{Language} 53(1):
    147--184.
  \item Goddard, Ives \& Bragdon, Kathleen. 1988. \emph{Native Writings in
    Massachusett}. Transactions of the American Philosophical Society, Vol.\ 78.
  \item O'Brien, Frank Waobu. \emph{Guide to Historical Spellings}. --- Master
    tables for spelling-to-phoneme mapping across SNEA languages.
  \item Salvucci, Claudio. \emph{InteractiveALR Vol.\ 27: Massachusetts}.
    Evolution Publishing.
  \item Trumbull, J. Hammond. 1903. \emph{Natick Dictionary}. Washington, D.C.:
    Bureau of American Ethnology.
  \item Williams, Roger. 1643. \emph{A Key into the Language of America}.
    London: Gregory Dexter.
  \item Winslow, Edward. 1624. \emph{Good Newes from New England}. London: I.D.
    for William Bladen and John Bellamie.
  \item Winslow, Edward. 2014. \emph{Good News from New England}. Ed.\ Kelly
    Wisecup. University of Massachusetts Press. ISBN 9781625340832.
\end{itemize}
